\documentclass[12pt]{article}
\usepackage[utf8]{inputenc}
\usepackage[T1]{fontenc}
\usepackage{hyperref}
\title{Wambon Language Resources for K-12 Curriculum Development}
\author{Compiled by Claude Sonnet 5.5 (claude-sonnet-5-5) for Andr\'as Kornai}
\date{}
\begin{document}
\maketitle

\section{Introduction}
Wambon (ISO 639-3 wms, Glottocode wamb1259) is a Trans-New Guinea language of Papua, Indonesia, in the Greater Awyu / Awyu-Dumut group. Speaker counts vary: Wikipedia gives about 3,900 and Joshua Project gives 5,100. The dialects are a northern one (Digul, including Ketum) and a southern one (Yonggom).

\section{Primary Educational Resources}
No primers or school materials were found.

\section{Linguistic and Reference Materials}
\begin{itemize}
\item Wikipedia article: \url{https://en.wikipedia.org/wiki/Wambon_language}
\item de Vries and de Vries-Wiersma (1992), \emph{The Morphology of Wambon of the Irian Jaya Upper-Digul Area}, KITLV Press. This was mentioned in search results, but no URL for it was seen.
\item Wilco van den Heuvel (2025), \emph{Yonggom Wambon, a Dumut language of West Papua}, an open-access volume that reanalyses Drabbe's 1959 description and includes a 500-item wordlist: \url{https://library.oapen.org/handle/20.500.12657/100430?show=full}
\end{itemize}
No standalone Wambon dictionary was found.

\section{Religious and Community Resources}
Joshua Project reports Bible portions (2013--2015) and a New Testament (2022--2023) in Wambon, with an online version on bible.is and a Jesus Film. The complete Bible does not exist. Joshua Project page: \url{https://joshuaproject.net/people_groups/15845}

\section{Digital and Audio Resources}
Joshua Project lists a Jesus Film and an online New Testament. Only the Joshua Project page above was seen in search results, so the individual items are not listed separately.

\section{Curriculum Development Context}
Wambon is not known to be taught in schools. The New Testament and the grammatical descriptions are the main written sources found. A curriculum would need community-led orthography and materials work. A language-based curriculum is not possible from these resources alone.

\end{document}